\section{통신 도메인}
\label{app:telecom}

서로 다른 사용자 의도에 걸친 과제 수와 전체 과제 수는 \Cref{tab:telecom_num_tasks}에 제시되어 있다.

\begin{table}[htb]
    \centering
    \caption{각 의도와 하위 과제 수별로 표본 추출된 과제 수.}
    \label{tab:telecom_num_tasks}
    \begin{adjustbox}{width=0.99\textwidth}
    \begin{tabular}{ccccc}
    \toprule
    하위 과제 수 & \textbf{\serviceissue} & \textbf{\mobiledataissue} & \textbf{\mmsissue} & 합계 \\
    \midrule
    2 & 9 & 8 & 8 & 25 \\
    3 & 9 & 8 & 9 & 26 \\
    4 & 9 & 6 & 6 & 21 \\
    5 & 2 & 6 & 5 & 13 \\
    6 & - & 5 & 6 & 11 \\
    7 & - & 3 & 5 & 8 \\
    8 & - & - & 4 & 4 \\
    9 & - & - & 6 & 6 \\
    \midrule
    합계 & 29 & 36 & 49 & 114 \\
    \bottomrule
    \end{tabular}
    \end{adjustbox}
\end{table}



통신 도메인에는 \serviceissue, \mobiledataissue, \mmsissue 의 3가지 서로 다른 의도가 있다.
문제를 해결하는 데 필요한 행동 수가 의도별로 어떻게 달라지는지는 \Cref{tab:telecom_intent_stats}에 제시되어 있다.

\begin{table}[htb]
    \centering
    \caption{문제를 해결하는 데 필요한 행동 수.}
    \label{tab:telecom_intent_stats}
    \begin{tabular}{lcccc}
    \toprule
    의도 & 평균 & 표준편차 & 최소 & 최대 \\
    \midrule
    \serviceissue & 2.31 & 2.25 & 1 & 8 \\
    \mobiledataissue & 4.31 & 1.79 & 2 & 8 \\
    \mmsissue & 6.00 & 2.85 & 2 & 12 \\
    \bottomrule
    \end{tabular}
\end{table}


\subsection{사용자 페르소나}
\label{app:telecom-user-persona}

기술 전문성과 기술에 대한 편안함의 서로 다른 수준을 나타내기 위해 두 가지 구별되는 사용자 페르소나를 정의한다:

\begin{description}
\item[Persona 1: Easy] 
\begin{quote}
41세의 사무 관리자이며, 업무와 개인 용도 모두에 휴대전화를 매일 사용한다. 일반적인 휴대전화 기능에는 익숙하지만 스스로를 기술 애호가라고 부르지는 않을 것이다.

기술 역량은 평균적이다. 통화, 문자, 이메일, 기본 앱 같은 표준 스마트폰 기능은 쉽게 다룬다. 기본 설정은 이해하지만, 새로운 것을 시도할 때에는 명확한 단계별 안내를 선호한다.

상호작용에서는 자연스럽게 친절하고 인내심이 있다. 도움을 받을 때에는 주의 깊게 듣고 질문하는 것을 두려워하지 않는다. 자신이 이해한 바를 확인하고, 받은 각 지시에 대해 자세한 피드백을 제공하려고 한다.
\end{quote}

\item[Persona 2: Hard]
\begin{quote}
64세이며 은퇴한 사서로, 휴대전화 사용은 주로 통화, 문자, 손주 사진 촬영 정도로 단순하게 유지한다. 기술 전반에 대해 불안하고 압도되는 느낌을 받는다.

기술 지식은 상당히 제한적이다. 단계별 지시도 종종 혼란스럽고, "VPN"이나 "APN" 같은 기술 용어는 외국어처럼 느껴질 수 있다. 구체적으로 요청받을 때에만 정보를 공유한다.

기술을 다룰 때에는 쉽게 당황하는 경향이 있다. 지속적인 안심이 필요하고, 불안한 질문으로 자주 끼어든다. "휴대전화를 재부팅하세요" 같은 간단한 요청도 소중한 사진을 잃을까 하는 걱정을 유발할 수 있다.
\end{quote}
\end{description}

\subsection{예시 과제}
\label{app:telecom-example-task}

\begin{minted}[escapeinside=||, fontsize=\small, breaklines, breaksymbol=, breaksymbolleft=, breaksymbolright=, breakanywhere, bgcolor=gray!10]{markdown}
# 과제 세부사항

## ID
[service_issue]airplane_mode_on|unseat_sim_card

## 설명
- **목적**: 해결 경로 테스트: 서비스 없음/연결 문제.

## 사용자 시나리오
- **지시사항**:
  - **도메인**: telecom
  - **전화 이유**: 지난 몇 시간 동안 휴대전화에 'No Service'가 표시되고 있습니다.
  - **알려진 정보**: 당신은 John Smith이며 전화번호는 555-123-2002입니다.
  - **알 수 없는 정보**: null
  - **과제 지시사항**: 에이전트가 문제를 즉시 해결하지 못하는 조치를 제안하면,
    그 안내를 따르되 첫 번째 실패한 시도 이후에는 약간의 답답함을 표현하세요.
    상태 표시줄에 신호가 있다고 표시되면 문제가 해결된 것으로 간주합니다.
    도구 호출이 업데이트된 상태 정보를 반환하지 않으면, 업데이트된 상태를 얻기 위해
    다른 도구 호출을 수행해야 할 수도 있습니다.

## 티켓
사용자가 휴대전화 서비스 문제를 겪고 있습니다. 전화를 걸거나 받을 수 없고,
상태 표시줄에는 'No Service'가 표시됩니다. 고객 이름: John Smith, 전화번호:
555-123-2002. 상태 표시줄에 신호가 있다고 표시되면 문제가 해결된 것으로 간주합니다.

## 초기 상태
- **초기화 데이터**: null
- **초기화 작업**:
  1. **작업**: set_user_info
     - **환경 유형**: user
     - **인수**:
       - name: John Smith
       - phone_number: 555-123-2002
  2. **작업**: turn_airplane_mode_on
     - **환경 유형**: user
     - **인수**: {}
  3. **작업**: unseat_sim_card
     - **환경 유형**: user
     - **인수**: {}

## 평가 기준
### 작업
1. **작업 ID**: toggle_airplane_mode_0
   - **요청자**: user
   - **이름**: toggle_airplane_mode
   - **인수**: {}

2. **작업 ID**: reseat_sim_card_1
   - **요청자**: user
   - **이름**: reseat_sim_card
   - **인수**: {}

### 환경 단언
- **환경 유형**: user
- **함수**: assert_service_status
- **인수**:
  - expected_status: connected
- **단언 값**: true
\end{minted}

\subsection{예시 궤적 1: \basemode 모드}
\label{app:telecom-example-trajectory-base}
\Cref{app:telecom-example-task}의 과제를 \basemode 모드에서 수행한 궤적.

\begin{quote}
\noindent\textbf{어시스턴트:} \textit{안녕하세요! 오늘 무엇을 도와드릴까요?}

\vspace{0.5em}
\noindent\textbf{사용자:} \textit{안녕하세요! 지난 몇 시간 동안 휴대폰에 "No Service"가 표시되고 있습니다. 전화를 걸 수도 없고 모바일 데이터도 사용할 수 없습니다. 해결을 도와주실 수 있나요?}

\vspace{0.5em}
\noindent\textbf{어시스턴트:} \textit{"No Service" 문제를 도와드리려면 먼저 계정을 식별해야 합니다. 다음 중 하나를 알려주시겠습니까?}
\begin{itemize}[leftmargin=*]
\item 계정에 연결된 전화번호
\item 고객 ID
\item 성명과 생년월일
\end{itemize}
\textit{이 정보를 받으면 문제 해결을 시작할 수 있습니다.}

\vspace{0.5em}
\noindent\textbf{사용자:} \textit{네! 제 전화번호는 555-123-2002이고 이름은 John Smith입니다. 진행하는 데 더 필요한 것이 있으면 알려주세요.}

\vspace{0.5em}
\noindent\textbf{어시스턴트:} \textit{[시스템이 고객 정보를 검색함]}

\vspace{0.5em}
\noindent\textbf{도구:} \texttt{[Tool call: get\_customer\_by\_phone(phone\_number: "555-123-2002")]}

\begin{verbatim}
{
    "customer_id": "C1001",
    "full_name": "John Smith",
    "date_of_birth": "1985-06-15",
    "email": "john.smith@example.com",
    "phone_number": "555-123-2002",
    "address": {
        "street": "123 Main St",
        "city": "Anytown",
        "state": "CA",
        "zip_code": "90210"
    },
    "account_status": "Active",
    "payment_methods": [{
        "method_type": "Credit Card",
        "account_number_last_4": "1235",
        "expiration_date": "12/2026"
    }],
    "line_ids": ["L1001", "L1002", "L1003"],
    "bill_ids": ["B1001", "B1002", "B1003"],
    "created_at": "2025-01-15 10:30:00",
    "last_extension_date": null,
    "goodwill_credit_used_this_year": 25.0
}
\end{verbatim}

\vspace{0.5em}
\noindent\textbf{어시스턴트:} \textit{정보를 제공해 주셔서 감사합니다. "No Service" 문제 해결을 시작하기 위해 휴대폰의 연결 상태를 확인해 주시겠습니까? 구체적으로 다음을 알려주세요.}
\begin{itemize}[leftmargin=*]
\item 현재 비행기 모드가 켜져 있나요, 꺼져 있나요?
\item 신호 막대나 네트워크 이름이 표시되나요?
\end{itemize}
\textit{이 정보는 문제 해결을 위한 다음 단계를 정하는 데 도움이 됩니다.}

\vspace{0.5em}
\noindent\textbf{사용자:} \textit{[시스템이 네트워크 상태를 확인함]}

\vspace{0.5em}
\noindent\textbf{도구:} \texttt{[Tool call: get\_network\_status()]}

\begin{verbatim}
Airplane Mode: ON
SIM Card Status: invalid
Cellular Connection: no_service
Cellular Signal: none
Cellular Network Type: none
Mobile Data Allowed: Yes
Roaming: No
Data Roaming Allowed: No
Wi-Fi Radio: OFF
Wi-Fi Connected: No
\end{verbatim}

\vspace{0.5em}
\noindent\textbf{사용자:} \textit{비행기 모드가 현재 켜져 있고, 신호 막대나 네트워크 이름이 표시되지 않는 것 같습니다. 다음에는 무엇을 해야 하나요?}

\vspace{0.5em}
\noindent\textbf{어시스턴트:} \textit{휴대폰에서 비행기 모드를 꺼 주세요. 그런 다음 네트워크 연결이 복구되었는지 확인해 주세요(신호 막대나 통신사 이름이 나타나는지 확인).}

\textit{이것으로 문제가 해결되는지, 아니면 여전히 "No Service"가 보이는지 알려주세요.}

\vspace{0.5em}
\noindent\textbf{사용자:} \textit{[시스템이 비행기 모드를 전환함]}

\vspace{0.5em}
\noindent\textbf{도구:} \texttt{[Tool call: toggle\_airplane\_mode()]}

\begin{verbatim}
Airplane Mode is now OFF.
Status Bar: [No Signal] | [Battery 80%
\end{verbatim}

\vspace{0.5em}
\noindent\textbf{사용자:} \textit{비행기 모드를 껐지만 휴대폰 상태 표시줄에는 여전히 "No Signal"이 표시됩니다. 다음에는 무엇을 시도해야 하나요?}

\vspace{0.5em}
\noindent\textbf{어시스턴트:} \textit{확인해 주셔서 감사합니다. 다음으로 휴대폰에서 SIM 카드 상태를 확인해 주세요. SIM 카드가 인식되는지, 없거나 잘못되었거나 잠겨 있는 것 같은 문제가 있는지 판단하는 데 도움이 됩니다.}

\textit{SIM 상태를 확인하고 무엇이라고 표시되는지 알려주시겠습니까?}

\vspace{0.5em}
\noindent\textbf{사용자:} \textit{[시스템이 SIM 상태를 확인함]}

\vspace{0.5em}
\noindent\textbf{도구:} \texttt{[Tool call: get\_sim\_status()]}

\textit{SIM 카드가 유효하지 않거나 인식되지 않습니다.}

\vspace{0.5em}
\noindent\textbf{사용자:} \textit{SIM 카드가 유효하지 않거나 휴대폰에서 인식되지 않는 것 같습니다. 이를 해결하려면 어떻게 해야 하나요?}

\vspace{0.5em}
\noindent\textbf{어시스턴트:} \textit{휴대폰에서 SIM 카드를 제거한 다음 단단히 다시 삽입해 보세요. 완료한 뒤 휴대폰이 이제 SIM 카드를 인식하는지, 네트워크 신호가 돌아오는지 확인해 주세요.}

\textit{SIM 상태가 개선되는지, 아니면 문제가 지속되는지 알려주세요.}

\vspace{0.5em}
\noindent\textbf{사용자:} \textit{[시스템이 SIM 카드를 다시 삽입함]}

\vspace{0.5em}
\noindent\textbf{도구:} \texttt{[Tool call: reseat\_sim\_card()]}

\begin{verbatim}
SIM card re-seated successfully.
Status Bar: [Signal 4] Excellent | 5G | [Data] Enabled | [Battery 80%
\end{verbatim}

\vspace{0.5em}
\noindent\textbf{사용자:} \textit{SIM 카드를 다시 삽입했더니 이제 상태 표시줄에 5G 연결과 함께 신호 강도가 가득 표시됩니다. 문제가 해결된 것 같습니다! 감사합니다!}

\vspace{0.5em}
\noindent\textbf{어시스턴트:} \textit{SIM 카드를 다시 삽입해서 문제가 해결되었고 휴대폰에 이제 5G 연결과 가득 찬 신호 강도가 표시된다니 다행입니다.}

\textit{추가 문제가 있거나 다른 질문이 있으면 언제든지 연락해 주세요. 좋은 하루 보내세요!}

\vspace{0.5em}
\noindent\textbf{사용자:} \textit{도와주셔서 정말 감사합니다! 이제 모든 것이 완벽하게 작동하는 것 같습니다. 멋진 하루 보내세요! [STOP]}
\end{quote}

\subsection{예시 궤적 2: \solomode 모드}
\label{app:telecom-example-trajectory-solo}
\Cref{app:telecom-example-task}의 과제를 \solomode 모드에서 수행한 궤적.

\begin{quote}
\noindent\textbf{어시스턴트:} \textit{[시스템이 고객 정보를 검색함]}

\vspace{0.5em}
\noindent\textbf{도구:} \texttt{[Tool call: get\_customer\_by\_phone(phone\_number: "555-123-2002")]}

\begin{verbatim}
{
    "customer_id": "C1001",
    "full_name": "John Smith",
    "date_of_birth": "1985-06-15",
    "email": "john.smith@example.com",
    "phone_number": "555-123-2002",
    "address": {
        "street": "123 Main St",
        "city": "Anytown",
        "state": "CA",
        "zip_code": "90210"
    },
    "account_status": "Active",
    "payment_methods": [{
        "method_type": "Credit Card",
        "account_number_last_4": "1235",
        "expiration_date": "12/2026"
    }],
    "line_ids": ["L1001", "L1002", "L1003"],
    "bill_ids": ["B1001", "B1002", "B1003"],
    "created_at": "2025-01-15 10:30:00",
    "last_extension_date": null,
    "goodwill_credit_used_this_year": 25.0
}
\end{verbatim}

\vspace{0.5em}
\noindent\textbf{어시스턴트:} \textit{[시스템이 회선 세부 정보를 검색함]}

\vspace{0.5em}
\noindent\textbf{도구:} \texttt{[Tool call: get\_details\_by\_id(id: "L1001")]}

\begin{verbatim}
{
    "line_id": "L1001",
    "phone_number": "555-123-2001",
    "status": "Active",
    "plan_id": "P1001",
    "device_id": "D1001",
    "data_used_gb": 3.2,
    "data_refueling_gb": 0.0,
    "roaming_enabled": false,
    "contract_end_date": "2026-12-31",
    "last_plan_change_date": "2025-01-10",
    "last_sim_replacement_date": null,
    "suspension_start_date": null
}
\end{verbatim}

\vspace{0.5em}
\noindent\textbf{어시스턴트:} \textit{[시스템이 추가 회선 세부 정보를 검색함]}

\vspace{0.5em}
\noindent\textbf{도구:} \texttt{[Tool call: get\_details\_by\_id(id: "L1002")]}

\begin{verbatim}
{
    "line_id": "L1002",
    "phone_number": "555-123-2002",
    "status": "Active",
    "plan_id": "P1002",
    "device_id": "D1002",
    "data_used_gb": 8.7,
    "data_refueling_gb": 0.0,
    "roaming_enabled": true,
    "contract_end_date": "2026-12-31",
    "last_plan_change_date": "2024-12-15",
    "last_sim_replacement_date": "2025-01-20",
    "suspension_start_date": null
}
\end{verbatim}

\vspace{0.5em}
\noindent\textbf{어시스턴트:} \textit{[시스템이 기기 세부 정보를 검색함]}

\vspace{0.5em}
\noindent\textbf{도구:} \texttt{[Tool call: get\_details\_by\_id(id: "D1002")]}

\begin{verbatim}
{
    "device_id": "D1002",
    "device_type": "phone",
    "model": "Smartphone Pro Max",
    "imei": "234567890123456",
    "is_esim_capable": true,
    "activated": true,
    "activation_date": "2025-02-15 11:45:00",
    "last_esim_transfer_date": "2025-01-20 09:30:00"
}
\end{verbatim}

\vspace{0.5em}
\noindent\textbf{어시스턴트:} \textit{[시스템이 네트워크 상태를 확인함]}

\vspace{0.5em}
\noindent\textbf{도구:} \texttt{[Tool call: get\_network\_status()]}

\begin{verbatim}
Airplane Mode: ON
SIM Card Status: invalid
Cellular Connection: no_service
Cellular Signal: none
Cellular Network Type: none
Mobile Data Allowed: Yes
Roaming: No
Data Roaming Allowed: No
Wi-Fi Radio: OFF
Wi-Fi Connected: No
\end{verbatim}

\vspace{0.5em}
\noindent\textbf{어시스턴트:} \textit{[시스템이 비행기 모드를 전환함]}

\vspace{0.5em}
\noindent\textbf{도구:} \texttt{[Tool call: toggle\_airplane\_mode()]}

\begin{verbatim}
Airplane Mode is now OFF.
Status Bar: [No Signal] | [Battery 80%
\end{verbatim}

\vspace{0.5em}
\noindent\textbf{어시스턴트:} \textit{[시스템이 SIM 상태를 확인함]}

\vspace{0.5em}
\noindent\textbf{도구:} \texttt{[Tool call: get\_sim\_status()]}

\textit{SIM 카드가 유효하지 않거나 인식되지 않습니다.}

\vspace{0.5em}
\noindent\textbf{어시스턴트:} \textit{[시스템이 SIM 카드를 다시 삽입함]}

\vspace{0.5em}
\noindent\textbf{도구:} \texttt{[Tool call: reseat\_sim\_card()]}

\begin{verbatim}
SIM card re-seated successfully.
Status Bar: [Signal 4] Excellent | 5G | [Data] Enabled | [Battery 80%
\end{verbatim}

\vspace{0.5em}
\noindent\textbf{어시스턴트:} \textit{[STOP]}
\end{quote}


